\section{Common readouts and economic policy qualification}\label{sec:compatible84}
An approximately optimized continuation is one function. Its coefficients determine the action in every observation cell. Separate action-switch tests deliberately forget part of this structure: an alternative may defeat the reference action without winning the complete feasible menu, and individually possible choices need not be jointly generated by the same readout. We retain these restrictions when transferring the original NBO qualification account to an economic decision. The attained continuation, economic primitives, action menu and policy-cost observations are unchanged.

\subsection{Winning actions and compatible choices}
Fix the finite acquisition partition $\mathcal C$, its original-law masses $p_C$, and a robust finite action menu $\mathcal A_C$. Each action in $\mathcal A_C$ is feasible throughout $C$. Let $\Theta=\prod_{j=1}^d[-B_j,B_j]$ be the declared readout box. With fitted labels and continuation policy $\tau$ held fixed, the implemented first-date rule minimizes
\[
 q_{Cb}(\theta)=c_{Cb}+v_{Cb}^{\mathsf T}\theta,
 \qquad b\in\mathcal A_C,
\]
using the original deterministic tie order. This is the fixed-feature own-future readout used in the preceding qualification construction, not an additional policy-value oracle. Write $F$ for its qualification objective, $\widehat\theta$ for the attained feasible readout, $U=F(\widehat\theta)$, and
\[
 \Theta_\epsilon=\{\theta\in\Theta:F(\theta)\le U+\epsilon\},
 \qquad \epsilon\ge0.
\]
The reference action in $C$ is $a_C$. The closed winner region
\begin{equation}\label{eq:winner84}
 W_{Cb}=\{\theta\in\Theta:q_{Cb}(\theta)\le q_{Ca}(\theta)
                     \text{ for every }a\in\mathcal A_C\}
\end{equation}
contains every readout that actually selects $b$. It can additionally contain ties resolved in favor of another action. Using this closed enlargement preserves the original tie rule and is conservative at its boundaries.

Suppose $D_{Cb}$ is a verified lower bound on $\inf_{W_{Cb}}F$. A separately certified empty winner region has bound $+\infty$. Retain the reference action and every alternative with $D_{Cb}\le U+\epsilon$ in $\mathcal A_C(\epsilon)$. An inherited reference-only switch bound remains valid and may be intersected with the full-menu bound by taking their maximum.

A joint event $H$ specifies actions $(C,b)$ at distinct cells. Suppose
\[
 D_H\le\inf\{F(\theta):\theta\in\bigcap_{(C,b)\in H}W_{Cb}\}
\]
is independently certified. If $D_H>U+\epsilon$, no member of $\Theta_\epsilon$ can select all the actions in $H$. Equality does not exclude the event. The collection of such verified events is denoted by $\mathcal H_\epsilon$; it need not contain every incompatible event.

Let $Q^\tau(x,b)$ be the actual original-law cost of taking $b$ and then following the common fixed continuation $\tau$. For each positive-mass cell, suppose numerical integration gives
\begin{equation}\label{eq:conditional-cost84}
 \ell_{Cb}\le
 \mathbb E[Q^\tau(X,b)-Q^\tau(X,a_C)\mid X\in C]
 \le u_{Cb}.
\end{equation}
Set $\ell_{Ca_C}=u_{Ca_C}=0$. Zero-mass cells contribute zero. All acquisition ambiguity, continuous-law integration and arithmetic allowances belong to these endpoints. The preceding actionwise calculation provides precisely these inputs.

\subsection{A computable economic bound}
For $z=(z_{Cb})$, define the assignment relaxation
\begin{equation}\label{eq:assignment84}
 \begin{split}
 \mathcal Z_\epsilon=\bigg\{z\ge0:\;&
 \sum_{b\in\mathcal A_C(\epsilon)}z_{Cb}=1\quad(C\in\mathcal C),\\
 &\sum_{(C,b)\in H}z_{Cb}\le |H|-1
                       \quad(H\in\mathcal H_\epsilon)\bigg\}.
 \end{split}
\end{equation}
Absent variables have value zero. Every policy induced by one readout in $\Theta_\epsilon$ defines a feasible zero--one assignment. Fractional assignments, and some unattainable zero--one assignments, can remain in $\mathcal Z_\epsilon$.

For any weights $w_{Cb}$, take numbers $\alpha_C\in\mathbb R$ and $\lambda_H\ge0$ satisfying
\begin{equation}\label{eq:assignment-dual84}
 \alpha_C+\sum_{H\ni(C,b)}\lambda_H\ge w_{Cb}
              \quad\bigl(b\in\mathcal A_C(\epsilon)\bigr),
\end{equation}
and define
\[
 V(w;\alpha,\lambda)=\sum_C\alpha_C+
                      \sum_{H\in\mathcal H_\epsilon}(|H|-1)\lambda_H.
\]
The dual need not be numerically optimal. Its feasibility and value can be checked exactly.

\begin{theorem}[Compatible optimizer-class cost bound]\label{thm:compatible84}
Under the preceding feasibility, winner, joint-event and conditional-cost conditions, let $V_\epsilon^+$ be any feasible dual value for $w_{Cb}=p_Cu_{Cb}$, and let $V_\epsilon^-$ be any feasible dual value for $w_{Cb}=-p_C\ell_{Cb}$. Then every $\theta\in\Theta_\epsilon$ satisfies
\begin{equation}\label{eq:compatible-cost84}
 -V_\epsilon^-\le J(\pi^\theta)-J(\pi^{\widehat\theta})
                    \le V_\epsilon^+.
\end{equation}
The exact assignment optima cannot give a wider interval than the full-menu cellwise extrema, which cannot give a wider interval than the inherited reference-only actionwise bound. Intersecting independently valid computed bounds preserves this ordering even when numerical dual proposals are not optimal.

For two constructed controllers $A,B$, suppose the attained-cost event gives
$L\le J(\pi^{\widehat\theta_A})-J(\pi^{\widehat\theta_B})\le U_0$.
Then all members of their respective qualified classes satisfy
\begin{equation}\label{eq:compatible-comparison84}
 L-V_A^- -V_B^+\le J(\pi^{\theta_A})-J(\pi^{\theta_B})
                       \le U_0+V_A^+ +V_B^-.
\end{equation}
The implication holds on that same event simultaneously for all allowances and all valid certificates, including certificates chosen after inspecting the covered intervals. It adds no failure probability for the continuum of readouts. A feasible numerical output with verified dual value $d\le\inf_\Theta F$ and $F(\theta)-d\le\epsilon$ belongs to $\Theta_\epsilon$.
\end{theorem}

\begin{proof}
Fix $\theta\in\Theta_\epsilon$. Its chosen action at each cell lies in the closed winner region. A singleton bound strictly larger than $U+\epsilon$ therefore cannot exclude that action. Similarly, selecting all members of a joint event in $\mathcal H_\epsilon$ would imply $F(\theta)\ge D_H>U+\epsilon$, a contradiction. Its choice indicators $z^\theta$ consequently belong to \eqref{eq:assignment84}.

The common continuation and the original initial law give the exact decomposition
\[
 J(\pi^\theta)-J(\pi^{\widehat\theta})
 =\sum_Cp_C\sum_bz^\theta_{Cb}
   \mathbb E[Q^\tau(X,b)-Q^\tau(X,a_C)\mid X\in C].
\]
Apply the upper inequality in \eqref{eq:conditional-cost84}. Multiplying \eqref{eq:assignment-dual84} by $z^\theta_{Cb}$, summing, and applying the nonnegative joint-event multipliers bounds the resulting expression by $V_\epsilon^+$. Apply the same argument to the negative lower endpoints for the other inequality in \eqref{eq:compatible-cost84}.

Removing all joint-event restrictions gives independent cell simplices. Their support is the sum of the cellwise maximal weights. Adding valid restrictions can only reduce that support. Full-menu exclusions retain a subset of the alternatives admitted by a reference-only comparison, proving the second ordering. Intersections preserve containment because each input bound contains the same cost difference.

Subtract the two instances of \eqref{eq:compatible-cost84} and add the attained-policy interval to obtain \eqref{eq:compatible-comparison84}. On the attained-cost event this entire argument is deterministic and holds for every valid collection of inequalities. Taking any number of allowances or any coefficient class therefore requires no additional probabilistic event. Finally,
$F(\theta)\le d+\epsilon\le\inf_\Theta F+\epsilon\le U+\epsilon$.
\end{proof}

A positive lower endpoint in \eqref{eq:compatible-comparison84} supports replacing $A$ by $B$ for every additional normalized charge below that endpoint. A strict interval contained in $(-m,m)$ gives equivalence at the declared margin $m$. These statements concern actual policy costs throughout the stated optimization classes. They do not identify a probability that an unrestricted fitting algorithm will enter those classes.

\subsection{Qualification with one common coefficient vector}
The fixed-label calibrated objective is
\[
 F(\theta)=\max_i(y_i^+-\phi_i^{\mathsf T}\theta)
           -\min_i(y_i^--\phi_i^{\mathsf T}\theta)
           +2\rho\bigl(G+\max_r|r_r^{\mathsf T}\theta|\bigr).
\]
A winner or joint-winner event adds linear inequalities
$d_j^{\mathsf T}\theta+c_j\le0$. The Numerical Appendix derives an exact lower certificate from two probability vectors, a derivative-budget multiplier and nonnegative event multipliers. All inequalities in a joint event use the same $\theta$ and the same qualification objective. A positive coefficient-box support separation certifies an empty event; a solver's infeasibility flag alone does not.

The complete-menu calculation and joint-event calculations can be stopped at any finite budget. Unchecked events add no exclusion, and inconclusive calculations retain the existing bound. Conversely, a valid new exclusion can immediately improve the economic assignment problem without any additional policy-cost observations. The arithmetic and integration costs required to obtain that exclusion remain separately chargeable.
